% Generated draft: requires mathematical and source review.
\documentclass[11pt]{article}
\usepackage[margin=25mm]{geometry}
\usepackage{fontspec}
\setmainfont{Latin Modern Roman}
\usepackage{amsmath,amssymb,mathtools,unicode-math}
\setmathfont{Latin Modern Math}
\usepackage{enumitem,needspace,xurl,hyperref}
\setlength{\parindent}{0pt}
\setlength{\parskip}{5pt}
\setlength{\emergencystretch}{3em}
\newcommand{\N}{\mathbb{N}}
\newcommand{\Z}{\mathbb{Z}}
\newcommand{\Q}{\mathbb{Q}}
\newcommand{\R}{\mathbb{R}}
\newcommand{\C}{\mathbb{C}}
\newcommand{\F}{\mathbb{F}}
\newcommand{\A}{\mathbb{A}}
\newcommand{\PP}{\mathbb{P}}
\newcommand{\E}{\mathbb{E}}
\newcommand{\eps}{\varepsilon}
\newcommand{\dd}{\,\mathrm{d}}
\begin{document}
\section*{UnsolvedMath ID: 30001399}
\textbf{Source problem code:} catalogue entry
\subsection{Definitions and mathematical statement}
For an operator \(T\), write \(\sigma(T)\) for its spectrum, once its underlying space and operator realization have been specified. The supplied source fragment uses \(h_{\eta,\lambda}\) for an operator associated with a configuration \(\eta\) and coupling parameter \(\lambda\), and \(\Sigma_\lambda\) for an almost-sure spectrum. It supplies neither the operator definition nor the probability model defining this almost-sure spectrum. In particular, it does not define the distinguished configurations \(\nu^*\) and \(\omega^1\).

The source asks whether, for every coupling parameter \(\lambda\) in the intended but unspecified parameter domain with \(|\lambda|>2\),
\[
\Sigma_\lambda
=
\sigma\!\left(h_{\nu^*,\lambda}\right)
\cup
\sigma\!\left(h_{\omega^1,\lambda}\right).
\]
This is the retained conjectural formula, not a self-contained mathematical statement: its interpretation requires the missing definitions of the model, probability law, operators, and configurations. The domain of \(\lambda\) must also be recovered; no real or complex domain is selected here.

The supplied literature assessment attributes a proof of the spectral formula for \(|\lambda|\leq 2\), including the boundary \(|\lambda|=2\), to the report, and describes the case \(|\lambda|>2\) as conjectural. These are supplied source claims and have not been independently verified. The registry labels the problem open but marks its statement \texttt{unrecoverable}, explaining that the undefined configurations prevent a publishable reconstruction.
\subsection{Short English statement}
The source conjectures that, for every intended coupling \(\lambda\) with \(|\lambda|>2\), the almost-sure spectrum equals the union of the spectra associated with \(\nu^*\) and \(\omega^1\). The configurations, operators, probability model, and parameter domain are missing, so this remains an incomplete statement requiring human mathematical review.
\subsection{Sources}
\begin{itemize}
\item Local UnsolvedMath registry, \texttt{data/unsolvedmath/problems.json}, permanent ID \(30001399\), title ``Incomplete Strong-Coupling Spectral Formula,'' exact source problem code \texttt{catalogue entry}. The supplied record belongs to the Oberwolfach Reports Open Problems collection, gives proposed year 2009, and names no individual proposer.
\item Supplied report citation: \emph{Mini-Workshop: Modeling and Understanding Random Hamiltonians: Beyond Monotonicity, Linearity and Independence} (2010), Oberwolfach Reports, pp. 2965--2992, DOI \texttt{10.4171/owr/2009/55}. Supplied source URL: \url{https://doi.org/10.4171/owr/2009/55}. The background describes Oberwolfach Reports as documenting mathematical workshops and their open-problem discussions.
\item Supplied literature source, labeled ``original Oberwolfach problem report'': \url{https://doi.org/10.4171/owr/2009/55}. Literature checked at: \texttt{2026-08-21}. This is supplied dated curation, not an independently verified assessment. It labels the problem open, attributes the spectral formula for \(|\lambda|\leq2\) to the report, and states that the report leaves \(|\lambda|>2\) conjectural and that no resolving primary result was located during that curation. The supplied background requests independent verification of time-sensitive status and cited claims.
\item Supplied statement-verification note: the extracted formula leaves \(\nu^*\) and \(\omega^1\) undefined; no separate clean statement is promoted, and the original record is retained publicly for provenance. The cited report is supplied as the primary reference, but its contents are not included beyond the extracted fragment and curation notes.
\end{itemize}
\end{document}
